\documentclass[10pt,a4paper]{article}
\usepackage[margin=0.7in]{geometry}
\usepackage{graphicx}
\usepackage{booktabs}
\usepackage{xcolor}
\usepackage{enumitem}
\usepackage{titlesec}
\usepackage{parskip}
\usepackage{amsmath}
\usepackage{hyperref}

\titleformat{\section}{\normalsize\bfseries\color{black!80}}{}{0pt}{}
\titlespacing*{\section}{0pt}{18pt}{6pt}
\setlength{\parskip}{2pt}
\pagestyle{empty}

\begin{document}

\begin{center}
  {\large\bfseries BTCUSDT Perpetual Futures: Microstructure Observations\\[2pt]
  and Live Paper-Trading Results}\\[4pt]
  {\footnotesize Mengren Man \quad$\cdot$\quad April 2026}
\end{center}

\vspace{6pt}\hrule\vspace{6pt}

%% ----------------------------------------------------------------
\section*{1.\ Overview}

Follow-up to the historical backtest summary shared on 9~April. This note
extends the original 2-day study in two directions:

\begin{enumerate}[leftmargin=1.5em, itemsep=1pt, topsep=2pt]
  \item \textbf{Tick-level microstructure analysis} on 7~days $\times$ 3~symbols
        (BTC, ETH, SOL USDT perpetual futures; 20M+ trades from
        \texttt{data.binance.vision}).
  \item \textbf{Live paper-trading} of the composite signal described in the
        original note, with a transparent taker cost model, running since
        10~April 2026.
\end{enumerate}

\noindent
No strategy changes have been made between the backtest and the live run.

%% ----------------------------------------------------------------
\section*{2.\ Tick-Level Microstructure Observations}

Data: 7~days of \texttt{aggTrades} for BTC/ETH/SOL USDT perpetuals
(2026-04-01 to 2026-04-07), totalling 20.3M trades and \$141B notional.

\subsection*{2.1\ Trade-sign autocorrelation (Lo--MacKinlay / Hasbrouck)}

In trade-time the aggressor sign is strongly persistent for the liquid names:

\smallskip
\begin{tabular}{lcccccc}
\toprule
 & $\rho_1$ & $\rho_5$ & $\rho_{10}$ & $\rho_{50}$ & $\rho_{100}$ & $\rho_{500}$ \\
\midrule
BTC & $+0.63$ & $+0.62$ & $+0.56$ & $+0.26$ & $+0.13$ & $+0.02$ \\
ETH & $+0.50$ & $+0.48$ & $+0.38$ & $+0.09$ & $+0.04$ & $+0.01$ \\
SOL & $-0.16$ & $+0.03$ & $+0.05$ & $+0.04$ & $+0.03$ & $+0.02$ \\
\bottomrule
\end{tabular}

\smallskip\noindent
A power-law fit $\rho_k \sim k^{-\gamma}$ on $k \in [5,\, 500]$ gives
$\gamma = 0.78$ (BTC) and $0.99$ (ETH), squarely within the equity-literature
range $[0.5,\, 1.0]$. The same order-splitting / metaorder signature that
governs equity tapes is present on crypto perpetual futures.

\textbf{SOL shows the opposite pattern}: $\rho_1 = -0.16$, with a shallow
decay $\gamma = 0.18$. This is the bid--ask-bounce signature of a book where
the 271\,ms mean inter-arrival is long enough that consecutive trades
alternate across the top of book. Two distinct microstructure regimes ---
order-splitting (BTC, ETH) and bid--ask bounce (SOL) --- are cleanly separated
by liquidity within the same 7-day sample.

\subsection*{2.2\ Trade-clock vs.\ wall-clock predictability}

AR(1) on signed flow in wall-clock buckets is uniformly weak:

\smallskip
\begin{tabular}{lccccc}
\toprule
 & 100\,ms & 500\,ms & 1\,s & 5\,s & 30\,s \\
\midrule
BTC $R^2$ & 0.008 & 0.003 & 0.005 & \textbf{0.023} & 0.016 \\
ETH $R^2$ & 0.004 & 0.003 & 0.007 & \textbf{0.034} & 0.018 \\
SOL $R^2$ & 0.001 & 0.001 & 0.001 & 0.005 & 0.012 \\
\bottomrule
\end{tabular}

\smallskip\noindent
Despite the large trade-time $\rho_1$ (0.50--0.63 for BTC/ETH), aggregating
into fixed time buckets destroys most of the sign signal.
\textbf{Trade-clock concentrates predictability; wall-clock dilutes it.}

\subsection*{2.3\ Additional observations}

Price-impact response is monotonic and roughly linear in signed size across
all three symbols, with concave tails on the largest prints (square-root-law
regime). Cross-symbol signed-flow lead--lag peaks at lag~0 at both 100\,ms and
10\,ms resolution --- no exploitable BTC$\to$alt lead on this sample. Realized
volatility scales with trade count, not wall-clock time (Clark 1973),
reinforcing the case for trade-clock or dollar-clock sampling in any
short-horizon signal.

%% ----------------------------------------------------------------
\section*{3.\ PCA Residual Mean-Reversion}

The original note's composite signal was extended with cross-asset PCA
residualization (BTC/ETH/SOL). PC1 explains 84\% of the joint variance;
stripping $\beta \cdot \text{PC1}$ from BTC yields a residual with lag-1
autocorrelation $\rho_1 = -0.11$ on 1-min time bars and $-0.13$ on \$5M
dollar bars (vs.\ $\approx 0$ for raw BTC returns).

\medskip
\begin{tabular}{llccc}
\toprule
\textbf{Experiment} & \textbf{Clock} & $\boldsymbol{\rho_1}$ & \textbf{Edge (bps/turn)} & \textbf{Net Sharpe @ 2\,bps} \\
\midrule
1-min raw BTC           & time  & $+0.00$ & ---   & --- \\
1-min PC1 residual      & time  & $-0.11$ & $\sim +0.2$ & $< 0$ \\
\$5M dollar bars, BTC only & dollar & --- & $-0.02$ & $< 0$ \\
\$5M dollar bars, PC1 residual & dollar & $-0.13$ & $+0.09$ & $\ll 0$ \\
\bottomrule
\end{tabular}

\medskip\noindent
The signal is real, robust, and cross-sample consistent --- but per-turn edge
is bounded above by $|\rho_1|\,\sigma^2 \approx 0.1$\,bps at the minute-ish
horizon. The constraint is execution infrastructure, not signal quality.

%% ----------------------------------------------------------------
\section*{4.\ Live Paper-Trading Results}

\textit{[This section will be populated with results from the live paper-trading
harness running since 10~April 2026.]}

\medskip\noindent
\textbf{Setup.} The composite signal from the original note (8 sub-signals,
default weights, VWAP reversion overweighted at 20\%, 5-bar EMA smoothing)
is recomputed on rolling 1-minute bars aggregated from a live websocket feed
(\texttt{btcusdt@aggTrade} + \texttt{btcusdt@bookTicker}).

\textbf{Fill model.} Taker-side: simulated market orders at live mid
$\pm$ half-spread, charged 2\,bps round-trip fee on $|\Delta\text{position}|$.
No queue-position modelling, no partial fills.

\medskip
\textit{[Insert: cumulative PnL chart, edge\_bps, live-vs-backtest tracking.
Purpose: pipeline sanity check and confirmation that the backtest was not
overfit --- not a claim of live profitability.]}

%% ----------------------------------------------------------------
\section*{5.\ What Would Close the Gap}

The per-turn edge of $\sim$0.1--0.2\,bps is structurally below any
realistic taker cost ($\sim$1--2\,bps round-trip). Three paths close the gap:

\begin{enumerate}[leftmargin=1.5em, itemsep=1pt, topsep=2pt]
  \item \textbf{Maker-side execution} --- fixes the \emph{cost sign}.
        Switch from crossing the spread to posting resting orders: collect
        spread $+$ exchange rebate instead of paying them. This converts the
        $\sim$1.5\,bps taker drag into a $\sim$1\,bps maker credit --- a
        2.5\,bps swing per turn that dominates the signal's own edge.
        Requires queue-position modelling and an exchange simulator to
        evaluate fill probability honestly.
  \item \textbf{Sub-second horizon} --- fixes the \emph{signal decay}.
        Trade-time sign autocorrelation ($\rho_1 = +0.63$ for BTC,
        section~2.1) is the directly tradable object at the sub-second
        horizon. $|\rho|$ is much larger there even though $\sigma^2$ per
        bar is smaller; the $|\rho|\,\sigma^2$ product can be
        10--100$\times$ the 1-minute value. This is the
        latency-sensitive regime where colocation and fast execution
        matter.
  \item \textbf{Broader cross-section} --- fixes the \emph{diversification}.
        Extending the PCA residual framework from 3~majors to 50+ altcoins
        makes individually marginal signals additive. A stacked portfolio
        of 15--20 decorrelated sub-bp signals can reach 1--3\,bps
        aggregate edge per turn.
\end{enumerate}

\noindent
All three are infrastructure or scale decisions, not research ones ---
and all three are core to how an HFT or systematic short-horizon firm
operates.

\end{document}
